\section{Ejemplos}


\begin{frame}{Aplicación: el problema de Basilea (I)}
    Sea $S_n(x)=\displaystyle\sum_{k=1}^n\frac{\sin kx}{k}$ en $[0,\pi]$. Por la teoría de series de Fourier,
    \[
        S_n(x)\to\frac{\pi-x}{2}\quad(0<x\le\pi),\qquad S_n(0)=0.
    \]
    El límite es discontinuo en $0$: la convergencia \textbf{no es uniforme}.

    \vspace{0.15cm}
    \textbf{Cota uniforme.} Sea $x\in(0,\pi]$ y $m=\lfloor1/x\rfloor+1$.
    \begin{itemize}
        \small
        \item Si $k<m$: $\abs{\sin kx}/k\le x$, y hay a lo sumo $1/x$ términos. Suman $\le1$.
        \item Si $k\ge m$: por la desigualdad de Abel, con $\frac1m<x$ y $\sin\frac{x}{2}\ge\frac{x}{\pi}$,
        \[
            \abs{\sum_{k=m}^n\frac{\sin kx}{k}}\le\frac1m\cdot\frac{1}{\sin(x/2)}\le x\cdot\frac{\pi}{x}=\pi.
        \]
    \end{itemize}
    Así, $\abs{S_n(x)}\le1+\pi$ para todo $n$ y todo $x\in[0,\pi]$.
\end{frame}

\begin{frame}{Aplicación: el problema de Basilea (II)}
    Por el teorema de Arzelà podemos integrar término a término:
    \[
        \sum_{k=1}^\infty\frac{1-(-1)^k}{k^2}
        =\lim_{n\to\infty}\int_0^\pi S_n(x)\,dx
        =\int_0^\pi\frac{\pi-x}{2}\,dx=\frac{\pi^2}{4}.
    \]
    Solo sobreviven los términos con $k$ impar, y cada uno aporta $2/k^2$:
    \[
        \sum_{k\text{ impar}}\frac{1}{k^2}=\frac{\pi^2}{8}.
    \]
    Separando pares e impares, $\zeta(2)=\dfrac{\pi^2}{8}+\dfrac{\zeta(2)}{4}$, y por tanto
    \[
        \boxed{\ \zeta(2)=\sum_{k=1}^\infty\frac{1}{k^2}=\frac{\pi^2}{6}\ }
    \]
\end{frame}

